%
%============= CHAPTER BREAK =================%
%
\chapter{Arrian of Nicomedia (108--160 CE)}
\vspace*{3mm}%
\justifying%
%

\noindent{}The introduction stated that the death of Hadrian %
would be the chronological point of departure of our neologistic exploration. %
Easier said than done. %
The adult life and writing career of our first author spanned the reigns of three emperors, %
Trajan, Hadrian, and Antoninus Pius. %
So while some of his works may be post-Hadrianic, %
many of them are not. %
The sorting out of the two requires careful consideration. %

Born in Nicomedia in Bithynia, %
Arrian traveled as a young man to Nicopolis of Epirus %
and became a student of Epictetus. % 
While the \textit{Discourses} or \textit{Dissertationes} of Epictetus %
are based on Arrianʼs notes from this time (107--108 CE), %
he did not publish this collection until after the death of Hadrian.\endnote{%
Syme 1982:209.} %
A given quotation from the \textit{Dissertationes}, then, %
could be considered as possibly reflective of Epictetus himself %
(in which case we plan to treat it in the second volume, currently in progress),\endnote{%
This approach dovetails with the retrospective methodology detailed and recommended in Vinzent 2019.} %
or as the product of Arrianʼs work as editor and publisher of the collection, %
in which case it appears in this chapter. %
Arrian was an enormously productive writer and well-known public figure, %
serving as legate and governor of Cappadocia from 131 to 137 CE, %
then elected an archon of Athens circa 145--147 CE,\endnote{%
Bobson 1976:x says 147 CE; Leon 2021:118.} and in the capital of Hellenistic culture he also died. \\[1mm] %

\noindent{}\textbf{οὐά} / \textit{oua} / “wah!” has 93 hits in the entire TLG. %
Some of these are inaccurate, textual corruptions, %
including the topmost TLG hit: Philodemus, \textit{On Rhetoric}.\endnote{Book 2Frag. col. z2.} %
Greek papyri and inscriptions have numerous matches for the three-letter string, %
but none appear to reflect the standalone exclamation. %
The lemma does appear three times in the \textit{Dissertationes} of %
Epictetus (3.22.34, 3.23.24, 3.23.32). %
It could thus be said to enter the literary record first either %
when Arrian took notes of Epictetusʼ dialogues ca 108 CE, %
or in the reign of Antoninus Pius when Arrian published them. %

In the NT, the lemma happens to be \textit{hapax legomenon}, occurring only in Mark 15:29, %
which recounts passersby during the crucifixion of Jesus who say, %
\textbf{Οὐὰ} ὁ καταλύων τὸν ναὸν καὶ οἰκοδομῶν ἐν τρισὶν ἡμέραις, %
/ “\textbf{Wah}, the one who destroys the temple and builds [it] in three days!” %
In my view, this unique feature and its use in this particular context of the crucifixion %
was likely inspired in part by Arrian: %

\vspace*{3mm}
 \begin{quote}εἰ γὰρ κακόν ἐστι τὸ ἀποθανεῖν, ἄν τε ὁμοῦ ἄν τε καθ’ ἕνα ὁμοίως κακόν ἐστιν. %
 μή τι ἄλλο τι μέλλει γίνεσθαι ἢ τὸ σωμάτιον χωρίζεσθαι καὶ ἡ ψυχή; οὐδέν. %
 (34) σοὶ δὲ ἀπολλυμένων τῶν Ἑλλήνων ἡ θύρα κέκλεισται; %
 οὐκ ἔξεστιν ἀποθανεῖν; ‘ἔξεστιν.’ τί οὖν πενθεῖς; %
 \textbf{‘οὐᾶ}, βασιλεὺς καὶ τὸ τοῦ Διὸς σκῆπτρον ἔχων.’ %
 ἀτυχὴς βασιλεὺς οὐ γίνεται· οὐ μᾶλλον ἢ ἀτυχὴς θεός.  (\textit{Diss} 3.22.34) \\[1mm] %
 For if dying is evil, then it is evil whether it is collective or individual. %
 Is anything else bound to happen than that the meager body and the soul are separated? Nothing. %
 (34) Has the door of the Greeks who died been shut to you? %
 Is it not permissible to die? “It is permissible.” So why do you mourn, %
 “\textbf{Wah}, a king and one who holds the scepter of Zeus!”? %
 A king does not become unfortunate, no more than a God [becomes] unfortunate. %
 \end{quote} % 
\vspace*{3mm}

\noindent{}The parallels go beyond the new, rare lemma, %
to include the foci of death, afterlife, monarchy, and divinity. %
Both also occur as the narratorʼs rendition of direct speech, %
as the first word of that speech act, %
and in specific relation to the death of a king. %
This imitation, if valid, suggests the Gospel of Mark (at least in its canonical form) %
was still being edited after 108 CE, and more likely late during, or subsequent to, the reign of Antoninus Pius, %
when Arrian published the \textit{Dissertationes}. %
That the next two witnesses to this lemma are Irenaeus (\textit{AH} 1.13.2) and Cassius Dio (\textit{H.R.} 63.20.5) %
in the late 2nd and early 3rd century favors a later time frame. %

One could argue that this lemma owes to Epictetus himself, %
or even to oral custom or a lost written source prior to Epictetus.\endnote{%
See Carrier 2026, who follows the outdated stenographic model of W.A. Oldfather %
when describing Arrian as one who merely endeavored “to transcribe” %
the street preaching of Epictetus. %
Carrier also claims it was “commonly used as an idiom” by Epictetus, %
one that reflected “the peopleʼs tongue” and was thus “clearly common dialect at the time”, %
i.e., “in the 50s and 60s AD”! These imaginative leaps, going well beyond %
any reasonable interpretation of the evidence in Arrianʼs \textit{Dissertationes} or any other extant Greek text, %
apparently serve to reinforce Carrierʼs fundamentalist date and depiction of %
canonical Mark as an integral product of a singular author of the 70s CE that was stabilized and maintained intact prior to its first manuscript witnesses.} %
The latter is an argument from silence, %
and the occurrence of this lemma in a text just three times, %
all within a relatively brief span, %
does not support the position that the lemma was commonplace custom, %
or reliably traced back to Epictetus as its source. %
Scholars continue to debate about the extent to which the \textit{Discourses} %
reflect the authentic teaching and voice of Epictetus the teacher %
vs. Arrian as his student, note-taker, editor, and compiler. %
Wirth has argued that it reflects Arrianʼs voice as a self-styled %
new Xenophon, with Epictetus implicitly portrayed as a new Socrates.\endnote{%
Wirth 1967.} %
On the opposite end of the spectrum, Dobbin has contended for %
Epictetus himself as author.\endnote{%
Dobbin 1988.} %
Several other scholars have subjected one or both of these extreme positions %
to scrutiny.\endnote{Brunt 1977; Stadter 1980; Stephens 1999; Long 2002:64.} %
We may safely say that Arrian took notes of the oral teachings of Epictetus, %
notes that he called “remembrances”, %
and later produced a stylistically distinctive %
and well-known compilation of these teachings. %
At the same time, Arrianʼs editorial voice, framing, and concerns %
are evident throughout. %
Thus the farthest back we can reliably trace this lemma is Arrian himself, %
not the decades-prior, spoken lectures or aphorisms of his teacher, %
still less any hypothetical sources or otherwise unattested, popular customs on which Epictetus supposedly drew. %